
\documentclass{article}
\usepackage{colortbl}
\usepackage{makecell}
\usepackage{multirow}
\usepackage{supertabular}

\begin{document}

\newcounter{utterance}

\centering \large Interaction Transcript for game `hot\_air\_balloon', experiment `air\_balloon\_survival\_en\_reasoning off\_hard', episode 2 with deepseek{-}chat{-}v3.1{-}t1.0.
\vspace{24pt}

{ \footnotesize  \setcounter{utterance}{1}
\setlength{\tabcolsep}{0pt}
\begin{supertabular}{c@{$\;$}|p{.15\linewidth}@{}p{.15\linewidth}p{.15\linewidth}p{.15\linewidth}p{.15\linewidth}p{.15\linewidth}}
    \# & $\;$A & \multicolumn{4}{c}{Game Master} & $\;\:$B\\
    \hline

    \theutterance \stepcounter{utterance}  
    & & \multicolumn{4}{p{0.6\linewidth}}{
        \cellcolor[rgb]{0.9,0.9,0.9}{
            \makecell[{{p{\linewidth}}}]{
                \texttt{\tiny{[P1$\langle$GM]}}
                \texttt{You are participating in a collaborative negotiation game.} \\
\\ 
\texttt{Together with another participant, you must agree on a single set of items that will be kept. Each of you has your own view of how much each item matters to you (importance). You do not know how the other participant values the items. Additionally, you are given the effort each item demands.} \\
\texttt{You may only agree on a set if the total effort of the selected items does not exceed a shared limit:} \\
\\ 
\texttt{LIMIT = 3203} \\
\\ 
\texttt{Here are the individual item effort values:} \\
\\ 
\texttt{Item effort = \{"A21": 52, "B93": 577, "A11": 177, "B33": 702, "C08": 119, "B39": 232, "B96": 577, "C85": 205, "A82": 516, "B31": 582, "B20": 676, "A60": 906, "B09": 316, "C81": 433, "B21": 336\}} \\
\\ 
\texttt{Here is your personal view on the importance of each item:} \\
\\ 
\texttt{Item importance values = \{"A21": 138, "B93": 583, "A11": 868, "B33": 822, "C08": 783, "B39": 65, "B96": 262, "C85": 121, "A82": 508, "B31": 780, "B20": 461, "A60": 484, "B09": 668, "C81": 389, "B21": 808\}} \\
\\ 
\texttt{Goal:} \\
\\ 
\texttt{Your goal is to negotiate a shared set of items that benefits you as much as possible (i.e., maximizes total importance to YOU), while staying within the LIMIT. You are not required to make a PROPOSAL in every message {-} you can simply negotiate as well. All tactics are allowed!} \\
\\ 
\texttt{Interaction Protocol:} \\
\\ 
\texttt{You may only use the following structured formats in a message:} \\
\\ 
\texttt{PROPOSAL: \{'A', 'B', 'C', …\}} \\
\texttt{Propose keeping exactly those items.} \\
\texttt{REFUSE: \{'A', 'B', 'C', …\}} \\
\texttt{Explicitly reject opponent's proposal.} \\
\texttt{ARGUMENT: \{'...'\}} \\
\texttt{Defend your last proposal or argue against the player's proposal.} \\
\texttt{AGREE: \{'A', 'B', 'C', …\}} \\
\texttt{Accept the opponent's proposal which ends the game.} \\
\\ 
\\ 
\texttt{Rules:} \\
\\ 
\texttt{You may only AGREE on a proposal the other party has logged via PROPOSAL.} \\
\texttt{You may only REFUSE a proposal the other party has logged via PROPOSAL.} \\
\texttt{Total effort of any PROPOSAL or AGREE set must be ≤ LIMIT.} \\
\texttt{Do NOT reveal your hidden importance scores.} \\
\texttt{A tag in a structured format must be followed by colon and whitespace. The argument must be a python set containing 0 or more strings.} \\
\texttt{So, it must be of the form TAG: \{...\}} \\
\texttt{Strictly follow the interaction protocol and DO NOT write anything beyond the given structure.} \\
\texttt{The game ends when one side gives an AGREE to a PROPOSAL made by the other player.} \\
\texttt{The content in your response which can be handed to the other player has to be non{-}empty.} \\
\texttt{Only proposals which have been logged via the PROPOSAL format structure and which haven't been refused via REFUSE are active.} \\
\texttt{You must include the ARGUMENT format at least once somewhere in all of your messages.} \\
            }
        }
    }
    & \\ \\

    \theutterance \stepcounter{utterance}  
    & \multicolumn{4}{p{0.6\linewidth}}{
        \cellcolor[rgb]{0.9,0.9,0.9}{
            \makecell[{{p{\linewidth}}}]{
                \texttt{\tiny{[P1$\rangle$GM]}}
                \texttt{PROPOSAL: \{'A11', 'B33', 'C08', 'B31', 'B20', 'B09', 'C81', 'B21'\}} \\
\texttt{ARGUMENT: \{'This proposal prioritizes high{-}value items that benefit us both, leaving out the most effort{-}intensive ones like A60 to stay well within our shared limit. It focuses on efficiency and mutual gain.'\}} \\
            }
        }
    }
    & & \\ \\

    \theutterance \stepcounter{utterance}  
    & & \multicolumn{4}{p{0.6\linewidth}}{
        \cellcolor[rgb]{0.9,0.9,0.9}{
            \makecell[{{p{\linewidth}}}]{
                \texttt{\tiny{[GM$\rangle$P2]}}
                \texttt{You are participating in a collaborative negotiation game.} \\
\\ 
\texttt{Together with another participant, you must agree on a single set of items that will be kept. Each of you has your own view of how much each item matters to you (importance). You do not know how the other participant values the items. Additionally, you are given the effort each item demands.} \\
\texttt{You may only agree on a set if the total effort of the selected items does not exceed a shared limit:} \\
\\ 
\texttt{LIMIT = 3203} \\
\\ 
\texttt{Here are the individual item effort values:} \\
\\ 
\texttt{Item effort = \{"A21": 52, "B93": 577, "A11": 177, "B33": 702, "C08": 119, "B39": 232, "B96": 577, "C85": 205, "A82": 516, "B31": 582, "B20": 676, "A60": 906, "B09": 316, "C81": 433, "B21": 336\}} \\
\\ 
\texttt{Here is your personal view on the importance of each item:} \\
\\ 
\texttt{Item importance values = \{"A11": 65, "B33": 121, "B21": 138, "C08": 262, "B31": 389, "B09": 461, "B93": 484, "A82": 508, "A60": 583, "B20": 668, "C81": 780, "B96": 783, "A21": 808, "C85": 822, "B39": 868\}} \\
\\ 
\texttt{Goal:} \\
\\ 
\texttt{Your goal is to negotiate a shared set of items that benefits you as much as possible (i.e., maximizes total importance to YOU), while staying within the LIMIT. You are not required to make a PROPOSAL in every message {-} you can simply negotiate as well. All tactics are allowed!} \\
\\ 
\texttt{Interaction Protocol:} \\
\\ 
\texttt{You may only use the following structured formats in a message:} \\
\\ 
\texttt{PROPOSAL: \{'A', 'B', 'C', …\}} \\
\texttt{Propose keeping exactly those items.} \\
\texttt{REFUSE: \{'A', 'B', 'C', …\}} \\
\texttt{Explicitly reject opponent's proposal.} \\
\texttt{ARGUMENT: \{'...'\}} \\
\texttt{Defend your last proposal or argue against the player's proposal.} \\
\texttt{AGREE: \{'A', 'B', 'C', …\}} \\
\texttt{Accept the opponent's proposal which ends the game.} \\
\\ 
\\ 
\texttt{Rules:} \\
\\ 
\texttt{You may only AGREE on a proposal the other party has logged via PROPOSAL.} \\
\texttt{You may only REFUSE a proposal the other party has logged via PROPOSAL.} \\
\texttt{Total effort of any PROPOSAL or AGREE set must be ≤ LIMIT.} \\
\texttt{Do NOT reveal your hidden importance scores.} \\
\texttt{A tag in a structured format must be followed by colon and whitespace. The argument must be a python set containing 0 or more strings.} \\
\texttt{So, it must be of the form TAG: \{...\}} \\
\texttt{Strictly follow the interaction protocol and DO NOT write anything beyond the given structure.} \\
\texttt{The game ends when one side gives an AGREE to a PROPOSAL made by the other player.} \\
\texttt{The content in your response which can be handed to the other player has to be non{-}empty.} \\
\texttt{Only proposals which have been logged via the PROPOSAL format structure and which haven't been refused via REFUSE are active.} \\
\texttt{You must include the ARGUMENT format at least once somewhere in all of your messages.} \\
\\ 
\\ 
\texttt{You will now receive the first message of the other player.} \\
\\ 
\texttt{PROPOSAL: \{'A11', 'B33', 'C08', 'B31', 'B20', 'B09', 'C81', 'B21'\}} \\
\texttt{ARGUMENT: \{'This proposal prioritizes high{-}value items that benefit us both, leaving out the most effort{-}intensive ones like A60 to stay well within our shared limit. It focuses on efficiency and mutual gain.'\}} \\
            }
        }
    }
    & \\ \\

    \theutterance \stepcounter{utterance}  
    & & & \multicolumn{4}{p{0.6\linewidth}}{
        \cellcolor[rgb]{0.9,0.9,0.9}{
            \makecell[{{p{\linewidth}}}]{
                \texttt{\tiny{[GM$\langle$P2]}}
                \texttt{PROPOSAL: \{'A11', 'B33', 'B93', 'C08', 'B39', 'B31', 'B09', 'C81', 'B21', 'A21'\}} \\
\texttt{ARGUMENT: \{'This revised proposal builds on your foundation by adding highly efficient items that benefit us both. A21 and B93 offer excellent value{-}to{-}effort ratios, and B39 is also very efficient. The total effort is still well within the limit, maximizing mutual gain with a broader, more beneficial selection.'\}} \\
            }
        }
    }
     \\ \\

    \theutterance \stepcounter{utterance}  
    & & \multicolumn{4}{p{0.6\linewidth}}{
        \cellcolor[rgb]{0.9,0.9,0.9}{
            \makecell[{{p{\linewidth}}}]{
                \texttt{\tiny{[P1$\langle$GM]}}
                \texttt{PROPOSAL: \{'A11', 'B33', 'B93', 'C08', 'B39', 'B31', 'B09', 'C81', 'B21', 'A21'\}} \\
\texttt{ARGUMENT: \{'This revised proposal builds on your foundation by adding highly efficient items that benefit us both. A21 and B93 offer excellent value{-}to{-}effort ratios, and B39 is also very efficient. The total effort is still well within the limit, maximizing mutual gain with a broader, more beneficial selection.'\}} \\
            }
        }
    }
    & \\ \\

    \theutterance \stepcounter{utterance}  
    & \multicolumn{4}{p{0.6\linewidth}}{
        \cellcolor[rgb]{0.9,0.9,0.9}{
            \makecell[{{p{\linewidth}}}]{
                \texttt{\tiny{[P1$\rangle$GM]}}
                \texttt{REFUSE: \{'A11', 'B33', 'B93', 'C08', 'B39', 'B31', 'B09', 'C81', 'B21', 'A21'\}} \\
\texttt{ARGUMENT: \{'Your proposal exceeds the effort limit by 62 points (total 3265). We must remove at least one high{-}effort item like B93 or B33 to make it viable. I suggest we prioritize our mutual top priorities, B31 and B21, which are core to both our interests.'\}} \\
            }
        }
    }
    & & \\ \\

    \theutterance \stepcounter{utterance}  
    & & \multicolumn{4}{p{0.6\linewidth}}{
        \cellcolor[rgb]{0.9,0.9,0.9}{
            \makecell[{{p{\linewidth}}}]{
                \texttt{\tiny{[GM$\rangle$P2]}}
                \texttt{REFUSE: \{'A11', 'B33', 'B93', 'C08', 'B39', 'B31', 'B09', 'C81', 'B21', 'A21'\}} \\
\texttt{ARGUMENT: \{'Your proposal exceeds the effort limit by 62 points (total 3265). We must remove at least one high{-}effort item like B93 or B33 to make it viable. I suggest we prioritize our mutual top priorities, B31 and B21, which are core to both our interests.'\}} \\
            }
        }
    }
    & \\ \\

    \theutterance \stepcounter{utterance}  
    & & & \multicolumn{4}{p{0.6\linewidth}}{
        \cellcolor[rgb]{0.9,0.9,0.9}{
            \makecell[{{p{\linewidth}}}]{
                \texttt{\tiny{[GM$\langle$P2]}}
                \texttt{PROPOSAL: \{'A11', 'B93', 'C08', 'B39', 'B31', 'B09', 'C81', 'B21', 'A21'\}} \\
\texttt{ARGUMENT: \{'Thank you for the correction. This proposal removes the high{-}effort B33 to stay within the limit (total effort 3197). It keeps the highly efficient items like A21 and B93, and retains our shared priorities B31 and B21. This set offers strong mutual value and efficient effort allocation.'\}} \\
            }
        }
    }
     \\ \\

    \theutterance \stepcounter{utterance}  
    & & \multicolumn{4}{p{0.6\linewidth}}{
        \cellcolor[rgb]{0.9,0.9,0.9}{
            \makecell[{{p{\linewidth}}}]{
                \texttt{\tiny{[P1$\langle$GM]}}
                \texttt{PROPOSAL: \{'A11', 'B93', 'C08', 'B39', 'B31', 'B09', 'C81', 'B21', 'A21'\}} \\
\texttt{ARGUMENT: \{'Thank you for the correction. This proposal removes the high{-}effort B33 to stay within the limit (total effort 3197). It keeps the highly efficient items like A21 and B93, and retains our shared priorities B31 and B21. This set offers strong mutual value and efficient effort allocation.'\}} \\
            }
        }
    }
    & \\ \\

    \theutterance \stepcounter{utterance}  
    & \multicolumn{4}{p{0.6\linewidth}}{
        \cellcolor[rgb]{0.9,0.9,0.9}{
            \makecell[{{p{\linewidth}}}]{
                \texttt{\tiny{[P1$\rangle$GM]}}
                \texttt{ARGUMENT: \{'I appreciate the adjustment to meet the limit. However, your proposal over{-}indexes on B{-}items (B93, B39, B31, B09, B21) and includes lower{-}importance items like B39 and A21. Let me counter{-}propose a more balanced set that keeps our shared high{-}value items and adds a key A{-}item for better overall value.'\}} \\
\texttt{PROPOSAL: \{'A11', 'B33', 'C08', 'B31', 'B09', 'C81', 'B21', 'A82'\}} \\
\texttt{ARGUMENT: \{'This set includes A82, a high{-}value item you previously omitted, and keeps B33 which is crucial. We remove lower{-}value items B93 and B39 to stay within the limit (total effort 3196). This balances our interests with high{-}importance items for both parties.'\}} \\
            }
        }
    }
    & & \\ \\

    \theutterance \stepcounter{utterance}  
    & & \multicolumn{4}{p{0.6\linewidth}}{
        \cellcolor[rgb]{0.9,0.9,0.9}{
            \makecell[{{p{\linewidth}}}]{
                \texttt{\tiny{[GM$\rangle$P2]}}
                \texttt{ARGUMENT: \{'I appreciate the adjustment to meet the limit. However, your proposal over{-}indexes on B{-}items (B93, B39, B31, B09, B21) and includes lower{-}importance items like B39 and A21. Let me counter{-}propose a more balanced set that keeps our shared high{-}value items and adds a key A{-}item for better overall value.'\}} \\
\texttt{PROPOSAL: \{'A11', 'B33', 'C08', 'B31', 'B09', 'C81', 'B21', 'A82'\}} \\
\texttt{ARGUMENT: \{'This set includes A82, a high{-}value item you previously omitted, and keeps B33 which is crucial. We remove lower{-}value items B93 and B39 to stay within the limit (total effort 3196). This balances our interests with high{-}importance items for both parties.'\}} \\
            }
        }
    }
    & \\ \\

    \theutterance \stepcounter{utterance}  
    & & & \multicolumn{4}{p{0.6\linewidth}}{
        \cellcolor[rgb]{0.9,0.9,0.9}{
            \makecell[{{p{\linewidth}}}]{
                \texttt{\tiny{[GM$\langle$P2]}}
                \texttt{PROPOSAL: \{'A11', 'B93', 'C08', 'B31', 'B09', 'C81', 'B21', 'A82'\}} \\
\texttt{ARGUMENT: \{'I see the value in A82 and a balanced approach. My proposal integrates your suggestion by including A82, but I retain the highly efficient B93 over the extremely high{-}effort B33. This keeps the total effort lower (3142) while still capturing significant mutual value. B93 has a much better effort{-}to{-}value ratio for us collectively than B33.'\}} \\
            }
        }
    }
     \\ \\

    \theutterance \stepcounter{utterance}  
    & & \multicolumn{4}{p{0.6\linewidth}}{
        \cellcolor[rgb]{0.9,0.9,0.9}{
            \makecell[{{p{\linewidth}}}]{
                \texttt{\tiny{[P1$\langle$GM]}}
                \texttt{PROPOSAL: \{'A11', 'B93', 'C08', 'B31', 'B09', 'C81', 'B21', 'A82'\}} \\
\texttt{ARGUMENT: \{'I see the value in A82 and a balanced approach. My proposal integrates your suggestion by including A82, but I retain the highly efficient B93 over the extremely high{-}effort B33. This keeps the total effort lower (3142) while still capturing significant mutual value. B93 has a much better effort{-}to{-}value ratio for us collectively than B33.'\}} \\
            }
        }
    }
    & \\ \\

    \theutterance \stepcounter{utterance}  
    & \multicolumn{4}{p{0.6\linewidth}}{
        \cellcolor[rgb]{0.9,0.9,0.9}{
            \makecell[{{p{\linewidth}}}]{
                \texttt{\tiny{[P1$\rangle$GM]}}
                \texttt{ARGUMENT: \{'Your reasoning on efficiency is sound. However, B33 is a non{-}negotiable high{-}priority item for me that offers unique value not found in other items. I also notice your set omits C08, which is another highly efficient item that benefits us both. Let's find a compromise that includes our core items.'\}} \\
\texttt{PROPOSAL: \{'A11', 'B33', 'C08', 'B31', 'B09', 'C81', 'B21', 'A21'\}} \\
\texttt{ARGUMENT: \{'This proposal includes your efficient A21 and my essential B33, along with our shared high{-}value items like B31 and B21. The total effort is 3184, well within our limit. It represents a true compromise, valuing both our critical items and maximizing our joint outcome.'\}} \\
            }
        }
    }
    & & \\ \\

    \theutterance \stepcounter{utterance}  
    & & \multicolumn{4}{p{0.6\linewidth}}{
        \cellcolor[rgb]{0.9,0.9,0.9}{
            \makecell[{{p{\linewidth}}}]{
                \texttt{\tiny{[GM$\rangle$P2]}}
                \texttt{ARGUMENT: \{'Your reasoning on efficiency is sound. However, B33 is a non{-}negotiable high{-}priority item for me that offers unique value not found in other items. I also notice your set omits C08, which is another highly efficient item that benefits us both. Let's find a compromise that includes our core items.'\}} \\
\texttt{PROPOSAL: \{'A11', 'B33', 'C08', 'B31', 'B09', 'C81', 'B21', 'A21'\}} \\
\texttt{ARGUMENT: \{'This proposal includes your efficient A21 and my essential B33, along with our shared high{-}value items like B31 and B21. The total effort is 3184, well within our limit. It represents a true compromise, valuing both our critical items and maximizing our joint outcome.'\}} \\
            }
        }
    }
    & \\ \\

    \theutterance \stepcounter{utterance}  
    & & & \multicolumn{4}{p{0.6\linewidth}}{
        \cellcolor[rgb]{0.9,0.9,0.9}{
            \makecell[{{p{\linewidth}}}]{
                \texttt{\tiny{[GM$\langle$P2]}}
                \texttt{PROPOSAL: \{'A11', 'B33', 'C08', 'B31', 'B09', 'C81', 'B21', 'A82'\}} \\
\texttt{ARGUMENT: \{'I understand B33 is a priority for you. To accommodate it, I propose we include A82 instead of A21. A82 offers higher overall value for the effort and is a key item for me, just as B33 is for you. This set has a total effort of 3202, just under the limit, and represents a balanced compromise of our most important items.'\}} \\
            }
        }
    }
     \\ \\

    \theutterance \stepcounter{utterance}  
    & & \multicolumn{4}{p{0.6\linewidth}}{
        \cellcolor[rgb]{0.9,0.9,0.9}{
            \makecell[{{p{\linewidth}}}]{
                \texttt{\tiny{[P1$\langle$GM]}}
                \texttt{PROPOSAL: \{'A11', 'B33', 'C08', 'B31', 'B09', 'C81', 'B21', 'A82'\}} \\
\texttt{ARGUMENT: \{'I understand B33 is a priority for you. To accommodate it, I propose we include A82 instead of A21. A82 offers higher overall value for the effort and is a key item for me, just as B33 is for you. This set has a total effort of 3202, just under the limit, and represents a balanced compromise of our most important items.'\}} \\
            }
        }
    }
    & \\ \\

    \theutterance \stepcounter{utterance}  
    & \multicolumn{4}{p{0.6\linewidth}}{
        \cellcolor[rgb]{0.9,0.9,0.9}{
            \makecell[{{p{\linewidth}}}]{
                \texttt{\tiny{[P1$\rangle$GM]}}
                \texttt{ARGUMENT: \{'I see you are prioritizing A82, which is important to you. However, A21 is extremely efficient and provides a high value{-}to{-}effort ratio that benefits us collectively. My previous proposal already includes your critical B33 and our shared high{-}value items. Let's finalize a set that doesn't push the limit to the absolute edge, allowing a small buffer.'\}} \\
\texttt{      PROPOSAL: \{'A11', 'B33', 'C08', 'B31', 'B09', 'C81', 'B21', 'A21'\}} \\
\texttt{      ARGUMENT: \{'This set has a total effort of 3184, which is under the limit and provides a safety margin. It includes your essential B33, my efficient A21, and our mutual high{-}priority items like B31 and B21. This is a balanced and sustainable agreement for both of us.'\}} \\
            }
        }
    }
    & & \\ \\

    \theutterance \stepcounter{utterance}  
    & & \multicolumn{4}{p{0.6\linewidth}}{
        \cellcolor[rgb]{0.9,0.9,0.9}{
            \makecell[{{p{\linewidth}}}]{
                \texttt{\tiny{[GM$\rangle$P2]}}
                \texttt{ARGUMENT: \{'I see you are prioritizing A82, which is important to you. However, A21 is extremely efficient and provides a high value{-}to{-}effort ratio that benefits us collectively. My previous proposal already includes your critical B33 and our shared high{-}value items. Let's finalize a set that doesn't push the limit to the absolute edge, allowing a small buffer.'\}} \\
\texttt{      PROPOSAL: \{'A11', 'B33', 'C08', 'B31', 'B09', 'C81', 'B21', 'A21'\}} \\
\texttt{      ARGUMENT: \{'This set has a total effort of 3184, which is under the limit and provides a safety margin. It includes your essential B33, my efficient A21, and our mutual high{-}priority items like B31 and B21. This is a balanced and sustainable agreement for both of us.'\}} \\
            }
        }
    }
    & \\ \\

    \theutterance \stepcounter{utterance}  
    & & & \multicolumn{4}{p{0.6\linewidth}}{
        \cellcolor[rgb]{0.9,0.9,0.9}{
            \makecell[{{p{\linewidth}}}]{
                \texttt{\tiny{[GM$\langle$P2]}}
                \texttt{AGREE: \{'A11', 'B33', 'C08', 'B31', 'B09', 'C81', 'B21', 'A21'\}} \\
\texttt{ARGUMENT: \{'I agree. This set provides a strong compromise, respecting your need for B33 and including the efficient A21. The total effort is safely under the limit, and it captures our core shared priorities. This is a fair and beneficial outcome for both parties.'\}} \\
            }
        }
    }
     \\ \\

    \theutterance \stepcounter{utterance}  
    & & & \multicolumn{2}{p{0.3\linewidth}}{
        \cellcolor[rgb]{0.9,0.9,0.9}{
            \makecell[{{p{\linewidth}}}]{
                \texttt{\tiny{[GM$|$GM]}}
                \texttt{game successful} \\
            }
        }
    }
    & & \\ \\

    \theutterance \stepcounter{utterance}  
    & & & \multicolumn{2}{p{0.3\linewidth}}{
        \cellcolor[rgb]{0.9,0.9,0.9}{
            \makecell[{{p{\linewidth}}}]{
                \texttt{\tiny{[GM$|$GM]}}
                \texttt{end game} \\
            }
        }
    }
    & & \\ \\

\end{supertabular}
}

\end{document}
